\section{Comparative Analysis}
\label{sec:comparison}

% Lead: Rajan. Code: docs/js/ (all variant engines), docs/trade-simulator.html, docs/process.html.
% Figures: generate from backtester data.

This section presents a cross-variant comparative analysis using the 26-season historical backtest, examining how each \cola{} variant ranks teams differently under identical historical data. We focus on three team-season case studies that isolate distinct properties of the mechanism family (sustained drought, circumvention via repeat tanking, and elimination-pool variance), then turn to the calibration of Capped \cola{}'s stockpile cap (Section~\ref{sec:case-cap-sweep}) and the handling of traded picks (Section~\ref{sec:trade-handling}, the largest source of implementation ambiguity identified in \citet{highley2026solvable4}).

\subsection{Case Study: Sacramento Kings (Maximum Drought)}
\label{sec:case-sac}

The Sacramento Kings accumulate the highest Classic \cola{} index in the dataset, driven by a 16-season non-playoff drought (2006--07 through 2021--22) interrupted only by a single first-round appearance in 2022--23. The Kings' Classic index reaches a peak of $11{,}250$ tickets in 2017--18 before diminishment from their \#2 pick that year partially resets the accumulation; they held the top-ranked Classic \cola{} lottery position for six consecutive seasons (2012--13 through 2017--18) and regained rank~1 in 2024--25 after the 2022--23 playoff interruption.

The three \cola{} variants reward this history through different channels:
\begin{itemize}
    \item \textbf{Classic \cola{}:} The flat $\alpha = 1{,}000$ annual increment compounds across missed-playoff seasons; the Kings' index grows monotonically until diminishment events (pick \#4 in 2021--22, pick \#2 in 2017--18) partially reset it. Ranking is driven almost entirely by accumulated carry-over.
    \item \textbf{Simple \cola{}:} Deterministic ranking by drought years places the Kings near the top of the 22-team eligible pool in most years, with win-tiebreak ordering shifting them to rank~2 in seasons where another non-playoff team has a slightly longer drought.
    \item \textbf{Countdown \cola{}:} McCarty number = drought $\times$ wins. Because the Kings typically won 25--33 games during their drought, their McCarty number is lower than the Classic \cola{} index would suggest; in 2024--25 they rank second overall (McCarty~280) behind Chicago (McCarty~390).
\end{itemize}

The Kings illustrate the core design philosophy of \cola{}: teams whose losing is \emph{organic} (genuine non-competitiveness across multiple seasons) receive the strongest reward. This is the case the mechanism is optimized for.

\subsection{Case Study: Philadelphia 76ers (Tanking Circumvention)}
\label{sec:case-phi}

The 2013--14 through 2016--17 Philadelphia 76ers (the ``Process'') represent the canonical strategic-tanking sequence in the modern NBA: four consecutive top-3 lottery outcomes (picks \#3, \#3, \#1, \#3) acquired through deliberate multi-season non-competitiveness. Under the current weighted lottery, each tanked season yielded a top-tier pick; under \cola{}, graduated diminishment makes the same sequence self-defeating.

\Cref{tab:phi-process} shows Philadelphia's Classic \cola{} trajectory under static backtest:

\begin{table}[t]
\centering
\small
\begin{tabular}{lccccc}
\toprule
Season & W--L & Real Pick & \cola{} Tickets & \cola{} Rank & Top Team (Tickets) \\
\midrule
2013--14 & 19--63 & \#3 & 2{,}750 & 8 / 14 & SAC (7{,}250) \\
2014--15 & 18--64 & \#3 & 2{,}375 & 9 / 14 & SAC (8{,}250) \\
2015--16 & 10--72 & \#1 & 2{,}188 & 10 / 14 & SAC (9{,}250) \\
2016--17 & 28--54 & \#3 & 1{,}000 & 14 / 14 & SAC (10{,}250) \\
2017--18 & 52--30 & --- & 375 & --- & SAC (11{,}250) \\
\bottomrule
\end{tabular}
\caption{Philadelphia 76ers under Classic \cola{} during ``The Process.'' Each top-3 pick triggers graduated diminishment (\#1: full reset; \#3: 50\% reduction), leaving Philadelphia dead last in lottery priority by the fourth tanking season despite a 10--72 record. The 2017--18 residual of 375 reflects a 25\% second-round playoff reduction applied to the post-draft 500 ticket balance (50\% of 1{,}000) after Philadelphia's second-round playoff appearance. Sacramento's monotonic climb in the right-hand column is the counterfactual the mechanism is designed to produce: organic non-competitiveness is rewarded cumulatively while the Process extracts diminishing returns.}
\label{tab:phi-process}
\end{table}

Two properties emerge:

\begin{enumerate}
    \item \textbf{Diminishing returns to repeat tanking.} The Classic \cola{} diminishment schedule ($\{1.0, 0.75, 0.5, 0.25\}$ fractional reductions for picks \#1--4) dominates the $\alpha = 1{,}000$ annual accumulation. A team that receives a top-3 pick in year $t$ cannot re-enter the top of the lottery in year $t+1$ through tanking alone; recovery requires several non-playoff seasons with no diminishment event.
    \item \textbf{Lottery-day attribution matters.} Philadelphia's 2017--18 index of 375 assumes Philadelphia absorbs the \#3 diminishment. Because Boston held the \#1 pick on lottery day (the Fultz--Tatum swap was post-lottery), Philadelphia is charged only for \#3, not \#1. Under draft-day attribution, Philadelphia would reset to zero, Boston would retain its full index, and the mechanism's reward structure would be inverted. \Cref{sec:lottery-day} documents the audit that established this correction for all 364 lottery picks in the dataset.
\end{enumerate}

\subsection{Case Study: Milwaukee Bucks (Elimination-Pool Variance)}
\label{sec:case-mil}

Countdown \cola{} produces probability distributions that diverge sharply from both Classic \cola{} and the status-quo weighted lottery. The Milwaukee Bucks in the 2024--25 backtest provide a compact illustration.

Milwaukee's McCarty number (drought $\times$ wins) is~144, placing them seventh overall in the 22-team Countdown pool---second in the second five-team pool, behind Memphis (also McCarty~144, tiebroken by wins) and trailing pool~1 leaders Chicago (McCarty~390), Sacramento (280), Los Angeles Clippers (200), Detroit (176), and Toronto (150). A naive reading of Milwaukee's position might anticipate a meaningful \#1-pick probability: they are in the top half of the pool, and share a McCarty number with the top-ranked team of their pool.

Monte Carlo simulation (10{,}000 trials) yields $\Prob[\text{pick } \#1 \mid \text{Milwaukee}] \approx 3.7\%$ and $\E[\text{pick} \mid \text{Milwaukee}] \approx 7.22$. For comparison, Chicago (the overall pool leader) has $\Prob[\text{pick } \#1] \approx 29.7\%$ and $\E[\text{pick}] \approx 2.52$, and Memphis (Milwaukee's pool-adjacent counterpart with the same McCarty) has $\Prob[\text{pick } \#1] \approx 5.1\%$ and $\E[\text{pick}] \approx 6.28$. Countdown \cola{} runs 21 sequential bottom-up draws across nested five-team pools, and the compounding of elimination probabilities concentrates \#1-pick mass heavily on pool-1 leaders while spreading the remainder across pool-2 teams in a pattern that amplifies the effect of intra-pool ticket differences. Analytic derivation of these probabilities requires tracking $2^{21}$ branch paths; the Monte Carlo approach (applied uniformly across all 26 seasons of the backtest) is the tractable alternative.

This case highlights a tradeoff: Countdown \cola{} satisfies Highley's ``no team falls more than 4 spots'' bounded-displacement property by design, but the same elimination structure introduces probability variance that weakens the correspondence between \cola{} index rank and expected draft position. The tradeoff is material for adoption: teams and owners evaluating the mechanism in advance cannot easily reason about expected outcomes without simulation.

\subsection{Capped \cola{} Cap Calibration}
\label{sec:case-cap-sweep}

Capped \cola{} (Definition~\ref{def:capped-cola}) is parameterised by a single tunable constant: the stockpile cap $\Max$. The cap controls the trade-off between (i) ordering resolution at the top of the lottery (a higher cap permits more separation between the highest-stockpile teams) and (ii) the per-series tanking exposure given by Lemma~\ref{lem:capped-series-bound} (a higher cap raises the maximum cost a team incurs by advancing past round~1). The empirical question is whether a value of $\Max$ exists that produces both ``no more than 3--4 teams at the cap simultaneously'' (preserves separation between teams in active rebuild) and ``per-series cost of advancing in the playoffs is small relative to a typical team's annual accumulation'' (defuses the playoff-tanking incentive the league flagged).

\Cref{tab:cap-sweep} reports the 21-stable-season aggregates from the sweep over $\Max \in \{75, 100, 125, 150, 175, 200\}$, excluding the cold-start window 2000--2004 (Section~\ref{sec:backtesting}).

\begin{table}[t]
\centering
\footnotesize
\begin{tabular}{rrrrrrr}
\toprule
$\Max$ & Teams at cap (mean / max) & Sep.\ R1--R5 (mean) & Per-series max & Per-series typical & Cum.\ regret (max) \\
\midrule
75  & 6.71 / 9 &  0.9 & 22.5 & 15 & 45.0 \\
100 & 4.33 / 8 & 11.7 & 30.0 & 20 & 53.7 \\
125 & 2.86 / 4 & 27.9 & 37.5 & 25 & 53.7 \\
\textbf{150} & \textbf{1.62 / 3} & \textbf{47.7} & \textbf{45.0} & \textbf{30} & \textbf{57.6} \\
175 & 0.95 / 2 & 65.4 & 52.5 & 35 & 67.2 \\
200 & 0.52 / 2 & 76.4 & 60.0 & 40 & 76.8 \\
\bottomrule
\end{tabular}
\caption{Sensitivity of Capped \cola{} aggregates to the cap value, computed across 21 stable seasons (2005--2025) of the 26-season backtest. Per-series ``max'' = $0.3 \cdot \Max$ (the play-in advancer round~1 case from Lemma~\ref{lem:capped-series-bound}); per-series ``typical'' = $0.2 \cdot \Max$ (any other series). Cumulative regret max is the worst empirical observation across the 21 seasons. The MAX = 150 row (bold) is the variant's published default per \citet{highley2026solvable5}.}
\label{tab:cap-sweep}
\end{table}

Three patterns emerge from the sweep:

\begin{enumerate}
    \item \textbf{Sub-100 caps over-compress the lottery.} At $\Max = 75$, more than six teams sit at the cap on average and the separation between the rank-1 and rank-5 stockpile drops to under 1 ticket. The mechanism's ordering signal is destroyed: teams with materially different long-run records are treated as interchangeable, defeating the priority-ordering property the variant inherits from the broader \cola{} family.
    \item \textbf{Caps above 175 erase the cap's purpose.} At $\Max \geq 175$, fewer than one team sits at the cap on average, meaning the cap rarely binds and the variant's per-series exposure converges to a fraction of an unbounded stockpile. The playoff-tanking-incentive argument the cap was introduced to make becomes uninformative.
    \item \textbf{$\Max = 150$ is on the frontier.} The cap-150 row produces the largest separation (47.7 tickets) compatible with the ``no more than 3--4 teams at cap'' target, while keeping the per-series max at 45 tickets (approximately a year and a half of accumulation for a team near the wins-increment ceiling). The closest alternative, $\Max = 125$, produces tighter exposure (37.5 ticket max per series) at the cost of compressing separation to 27.9 tickets and pushing the typical teams-at-cap count to 2.9. Our memo-stage recommendation was $\Max = 125$ \citep{rajan2026memo}; \citet{highley2026solvable5} selects $\Max = 150$ on the rationale that the wider separation is worth the additional 7.5 tickets of per-series exposure.
\end{enumerate}

The published $\Max = 150$ choice produces fewer than two teams at the cap in 19 of 21 stable seasons, and at most three in any single season. The cap binds primarily in the late 2010s, when accumulated wins for the league's longest-rebuilding franchises (Sacramento, Detroit, Charlotte) approach the cap simultaneously. \highley{We have read the $\Max = 150$ choice as a published default in \citet{highley2026solvable5}; the Substack post phrases it as the calibration that supports the headline ``no playoff series win ever costs a team more than 45 lottery tickets.'' If you would prefer the paper present $\Max$ as a free parameter and report results across the sweep without privileging 150, we can demote the bold row in Table~\ref{tab:cap-sweep} and reframe the section as a calibration exercise. Our recommendation is to keep 150 as the default and present 125 as the closest credible alternative, mirroring the memo.}

\subsection{Trade Handling}
\label{sec:trade-handling}

\citet{highley2026solvable4} identifies four options for handling traded picks under \cola{}:

\begin{description}
    \item[Option 1 (Exclude):] Traded picks do not affect either team's \cola{} index.
    \item[Option 2 (Original Owner):] Diminishment applies to the team that originally held the pick (i.e., the team whose record determined the lottery position).
    \item[Option 3 (Receiving Team):] Diminishment applies to the team that ultimately receives the pick via trade.
    \item[Option 4 (Split):] Diminishment is divided between the original owner and receiving team.
\end{description}

Our dataset contains 55 verified pre-draft pick trades across 26 seasons, of which 8 fall in the top-3 diminishment range where the choice of rule materially affects the index (pick~4 trades did not occur in our window). The BKN~2017 \#1 pick (traded to BOS via the 2013 Billy King mega-trade) illustrates the rule-dependent index trajectory most sharply: \cref{tab:trade-options} reports BKN's Classic \cola{} index both \emph{before} the 2017 lottery (which captures the cumulative effect of 2016's BKN $\to$ BOS \#3 trade under each rule) and \emph{after} 2017's \#1 diminishment fires.

\begin{table}[t]
\centering
\footnotesize
\begin{tabular}{lrrrr}
\toprule
Rule & BKN pre-2017 & BKN post-2017 & BOS pre-2017 & BOS post-2017 \\
\midrule
Opt.\ 1 (Exclude)         & 4{,}635 & 4{,}635 & 500 & 500 \\
Opt.\ 2 (Original owner)  & 2{,}346 &     0   & 500 & 500 \\
Opt.\ 3 (Receiving team)  & 4{,}635 & 4{,}635 & 250 &   0 \\
Opt.\ 4 (Split)           & 3{,}373 & 1{,}687 & 375 & 188 \\
\bottomrule
\end{tabular}
\caption{BKN $\to$ BOS \#1 pick in 2017 under each trade rule. ``Pre-2017'' shows the entering lottery-day index under the cumulative application of each rule through prior seasons (the $\approx\!2{,}300$ spread across rules is driven by 2016's BKN $\to$ BOS \#3 trade). ``Post-2017'' shows the post-draft index after 2017's \#1 diminishment. The pre-index range (2{,}346 to 4{,}635) matches \citet{highley2026solvable4}'s qualitative claim that pick-tracking choice materially shifts the manipulator's carried-forward state. Values computed by \texttt{scripts/audit\_paper\_numbers.js} against the 26-season backtest.}
\label{tab:trade-options}
\end{table}

The two failure modes are visible in the data:

\begin{itemize}
    \item \textbf{Over-penalization under Option 2.} The 2011 LAC $\to$ CLE \#1 pick (the Baron Davis trade, which sent an unprotected LAC first-rounder to Cleveland years before the 2010--11 season) penalizes LAC's carried-forward index for a pick they had traded away with no expectation it would be \#1---under Option 2, LAC's index drops to zero after the 2011 draft despite their 2010--11 record having no tanking intent tied to the pick outcome. LAC's tanking incentive is severed from the pick outcome.
    \item \textbf{Rewarding the wrong actor under Option 3.} The 2016 BKN $\to$ BOS \#3 case (from the same 2013 Billy King mega-trade) illustrates the converse: BKN tanked (though arguably due to organizational dysfunction rather than strategic intent) and BOS received the \#3 pick. Under Option~3, BOS's index is reduced for a pick they did not earn through their own record ($1{,}000 \to 500$), while BKN's large accumulated index ($3{,}635$) is unaffected.
\end{itemize}

Option 4 (split) distributes diminishment evenly but compounds the underlying attribution problem: neither the manipulator nor the beneficiary absorbs the full signal.

Option 1 (exclude) is the simplest but eliminates the diminishment mechanism's effect on any traded pick, reintroducing a manipulation channel (trade the pick, then tank).

\citet{highley2026solvable4} recommends Option 2 paired with a strengthened Stepien Rule (cap on the number of future firsts a team can trade), which limits the frequency with which Option 2's over-penalization can occur while preserving the linkage between manipulation incentive and diminishment.

\subsection{Cross-Variant Summary}
\label{sec:cross-variant}

\Cref{tab:cross-variant} summarises how the five variants treat the three case-study teams, the cap-calibration result, and the primary trade-handling concern.

\begin{table}[t]
\centering
\footnotesize
\begin{tabular}{lp{2.0cm}p{2.0cm}p{2.0cm}p{2.0cm}p{2.4cm}}
\toprule
Property & Simple & Simple Lot. & Classic & Countdown & Capped \\
\midrule
Top team 2024--25 & CHI & CHI & SAC & CHI & SAC (at cap) \\
Top team $\Prob[\#1]$ & N/A (det.) & $25.0\%$ & $14.3\%$ & $29.7\%$ & raffle by $L_i / P$ \\
PHI 2016--17 rank & 21/22 & 21/22 & 14/14 & 21/22 & ineligible (drought $<2$) \\
Lottery component & None & Pre-2019 odds & Weighted & Elim.\ pools (MC) & Top-5 raffle, then by rank \\
Bounded displacement & 0 (det.) & $\pm 13$ & $\pm 13$ & $\pm 4$ (design) & by stockpile rank \\
Per-series tanking cost & 1 drought yr & 1 drought yr & 0--$\rho \cdot L_i$ & McCarty-dep. & $\leq 0.3 \cdot \Max$ \\
Clean trade handling & N/A & N/A & Option-dep. & Option-dep. & Option 2 + Stepien \\
Fan explainability & High & High & Medium & Low & Medium-low \\
\bottomrule
\end{tabular}
\caption{Cross-variant comparison. Numeric entries are drawn from the 2024--25 backtest; PHI 2016--17 rank is shown relative to the eligible pool size for each variant (14 for Classic, 22 for Simple/Simple Lottery/Countdown; ineligible under Capped \cola{} because the 2016--17 76ers had received a top-5 lottery pick the prior season). ``Bounded displacement'' = the maximum rank-vs-pick gap the mechanism admits; Countdown \cola{} enforces $\pm 4$ by design via the five-team elimination-pool structure. ``Per-series tanking cost'' = the maximum reduction in priority a team incurs from advancing one playoff series; Capped \cola{}'s bound is from Lemma~\ref{lem:capped-series-bound}. SAC leads Classic \cola{} but CHI leads the drought-ranked variants: Classic's cumulative carry-over rewards Sacramento's 16-season drought; drought-based ranking in the other variants puts Chicago's uninterrupted recent stretch first. Under Capped \cola{} at $\Max = 150$, both teams sit at the cap (Section~\ref{sec:case-cap-sweep}).}
\label{tab:cross-variant}
\end{table}

Four observations follow:

\begin{enumerate}
    \item \textbf{Rank agreement at the bottom, disagreement at the top.} All five variants agree that the 2016--17 76ers should receive near-lowest lottery priority. Capped \cola{} agrees by a different mechanism (eligibility): the 76ers' 2015--16 \#1 pick removes them from the eligible pool entirely under the drought $\geq 2$ requirement, while the other four variants rank them at or near the bottom of an eligible pool. At the top of the 2024--25 pool, however, the variants disagree: drought-based ranking (Simple, Simple Lottery, Countdown) puts Chicago first, while Classic \cola{}'s cumulative ticket carry-over puts Sacramento first. Capped \cola{} renders the question moot in 2024--25 by placing both at the cap.
    \item \textbf{Lottery structure drives probability variance.} $\Prob[\#1]$ for the top-ranked team ranges from $14.3\%$ (Classic, diffuse weighted lottery) to $29.7\%$ (Countdown, concentrated pool-1 leader) on the same 2024--25 data. Simple \cola{} is deterministic; Simple Lottery applies fixed 25\% worst-team odds regardless of pool state; Capped \cola{} raffles the top five picks proportional to stockpile share with multiple teams at the cap effectively splitting the highest-priority slot. Owners evaluating the mechanism in advance cannot treat ``top rank'' as equivalent to ``high \#1 probability'' across variants.
    \item \textbf{Complexity--robustness tradeoff.} Countdown \cola{} bounds displacement tightly but requires Monte Carlo simulation to reason about probabilities. Simple \cola{} is fully transparent but loses the randomisation that protects against degenerate strategic behaviour at the ranking boundary. Classic \cola{} occupies the middle and is the most common basis for comparison with the status quo. Capped \cola{} adds a fifth point on the frontier: it preserves the rank-driven raffle structure of Classic but bounds the per-series tanking cost via the cap (Lemma~\ref{lem:capped-series-bound}), at the cost of reducing the resolution between teams above the cap.
    \item \textbf{Capped \cola{} occupies a distinct adoption slot.} The four pre-Capped variants form a frontier from full transparency (Simple) to bounded displacement (Countdown). Capped \cola{} sits off this frontier: it sacrifices some between-team resolution to bound the per-series cost in a way the others do not. This is the design feature \citet{highley2026solvable5} introduces in response to the league's ``flexible lottery line is not on the table'' feedback; it is the only variant that addresses the playoff-tanking incentive directly through a stockpile bound rather than through eligibility or schedule.
\end{enumerate}

These tradeoffs frame the adoption discussion in \cref{sec:discussion}: no single variant dominates the others on all axes, which is itself an argument for presenting \cola{} as a family of mechanisms rather than a single design.
